% Options for packages loaded elsewhere
\PassOptionsToPackage{unicode}{hyperref}
\PassOptionsToPackage{hyphens}{url}
\documentclass[
  10pt,
  a4paper,
]{article}
\usepackage{xcolor}
\usepackage[margin=20mm]{geometry}
\usepackage{amsmath,amssymb}
\setcounter{secnumdepth}{5}
\usepackage{iftex}
\ifPDFTeX
  \usepackage[T1]{fontenc}
  \usepackage[utf8]{inputenc}
  \usepackage{textcomp} % provide euro and other symbols
\else % if luatex or xetex
  \usepackage{unicode-math} % this also loads fontspec
  \defaultfontfeatures{Scale=MatchLowercase}
  \defaultfontfeatures[\rmfamily]{Ligatures=TeX,Scale=1}
\fi
\usepackage{lmodern}
\ifPDFTeX\else
  % xetex/luatex font selection
\fi
% Use upquote if available, for straight quotes in verbatim environments
\IfFileExists{upquote.sty}{\usepackage{upquote}}{}
\IfFileExists{microtype.sty}{% use microtype if available
  \usepackage[]{microtype}
  \UseMicrotypeSet[protrusion]{basicmath} % disable protrusion for tt fonts
}{}
\makeatletter
\@ifundefined{KOMAClassName}{% if non-KOMA class
  \IfFileExists{parskip.sty}{%
    \usepackage{parskip}
  }{% else
    \setlength{\parindent}{0pt}
    \setlength{\parskip}{6pt plus 2pt minus 1pt}}
}{% if KOMA class
  \KOMAoptions{parskip=half}}
\makeatother
\usepackage{longtable,booktabs,array}
\usepackage{caption}
\captionsetup[table]{skip=6pt}
\newcounter{none} % for unnumbered tables
\usepackage{calc} % for calculating minipage widths
% Correct order of tables after \paragraph or \subparagraph
\usepackage{etoolbox}
\makeatletter
\patchcmd\longtable{\par}{\if@noskipsec\mbox{}\fi\par}{}{}
\makeatother
\setlength{\emergencystretch}{3em} % prevent overfull lines
\providecommand{\tightlist}{%
  \setlength{\itemsep}{0pt}\setlength{\parskip}{0pt}}

\usepackage{amsmath,amssymb,amsthm}
\setlength{\emergencystretch}{2em}
\setlength{\parskip}{4pt}
\setlength{\parindent}{0pt}
\usepackage{bookmark}
\IfFileExists{xurl.sty}{\usepackage{xurl}}{} % add URL line breaks if available
\urlstyle{same}
% fallback for those not using the hyperref driver hyperxmp:
\makeatletter
\@ifundefined{xmpquote}{\newcommand{\xmpquote}[1]{#1}}{}
\makeatother
\hypersetup{
  pdftitle={On the touching arc of the snub cube: a closed form and an independent check},
  pdfauthor={Sumin Yoo},
  hidelinks,
  pdfcreator={LaTeX via pandoc}}

\title{On the touching arc of the snub cube: a closed form and an
independent check}
\author{Sumin Yoo\\{\normalsize Independent researcher, Republic of Korea}}
\date{Version 1.0 -- 2026-10-15\\{\small Note text and PDF: \href{https://creativecommons.org/licenses/by/4.0/}{CC BY 4.0}.}}

\begin{document}
\maketitle

The snub cube admits an initial arc of poses in which one congruent
shadow is contained in another with boundary contact. L, W established
this arc and enclosed its endpoint numerically in Theorem 4.2 of their
preprint. We give the compact expression \(\theta_*=\arctan(2-t)\),
where \(t\) is the tribonacci constant, and an independent certificate
using exact arithmetic in \(\mathbb Q(t)\) and Arb interval arithmetic.
A change of coordinates identifies this arc with the tie described in
Tom 7's repository; its angular parameter has the same endpoint, whereas
its Cayley parameter ends at \(2-t\). Thus the repository's approximate
value \(0.165\) overstates the initial touching range under either
interpretation. Two separate remarks concern a weight-sign inference and
a binary64 bound in the L, W text. The scope is this arc and these
textual points; no global non-Rupert certificate is reproduced here.

\section{Setup and statement}\label{setup-and-statement}

Let \(t\in(1,2)\) be the unique real root of \(t^3-t^2-t-1=0\), and put
\(q=1/t=t^2-t-1\). Let \(V\) comprise the even permutations of
\((\pm1,\pm q,\pm t)\) with an even number of plus signs and the odd
permutations with an odd number of plus signs. The 24-point convex hull
\(P=\operatorname{conv}V\) is one chirality of the snub cube. A uniform
scaling or a simultaneous orthogonal change of coordinates does not
affect the containment question.

Write

\[a_4=e_z,\quad C=(q-1,q+1,0),\quad B=C\times a_4=(q+1,1-q,0),\]

\[u_\theta=\widehat C\cos\theta-\widehat B\sin\theta,\qquad
K_n=I-\frac{2nn^{\mathsf T}}{n\cdot n},\qquad K_0=K_C.\]

Here a hat denotes normalization, and \(\pi_u\) denotes orthogonal
projection onto \(u^\perp\). Rotations are active, right-handed
rotations. In particular, the \(xz\)-block of \(R_y(\alpha)\) is
\(\left(\begin{smallmatrix}\cos\alpha&\sin\alpha\\-\sin\alpha&\cos\alpha\end{smallmatrix}\right)\).
We ask for \textbf{closed containment at zero translation}, not
containment in the interior. For unit \(u\), define the clearance of a
fixed pose by
\(\min_{d\perp u,\ |d|=1}(h_{\rm hole}(d)-h_{\rm plug}(d))\). Thus
closed containment gives nonnegative clearance, and contact gives zero
clearance.

\textbf{Theorem.} For \(0\leq\theta\leq\arctan(1/3)\),

\[\pi_{u_\theta}(K_0P)\subseteq\pi_{u_\theta}(P)
\quad\Longleftrightarrow\quad
0\leq\theta\leq\boxed{\arctan(2-t)}.\]

Equivalently the plug can be \(R_z(2\theta)P\), with no translation. The
clearance is exactly zero throughout this initial arc. Every angle
immediately beyond its endpoint, up to \(\arctan(1/3)\), fails closed
containment for this pose. The angular domain is explicit because the
family is periodic; this is not an assertion that containment occurs
only once among all real angles.

\newpage

\section{Proof and certificate}\label{proof-and-certificate}

Set \(s=\tan\theta\). Since \(|B|=|C|\), coordinates for the screen,
after a positive horizontal rescaling, are

\[Q_s(v)=(B\cdot v+sC\cdot v,v_z),\qquad
Q'_s(v)=(B\cdot v-sC\cdot v,v_z).\]

The latter are the coordinates of \(K_0v\). The rescaling is invertible
and hence preserves containment and the signs of oriented areas. In a
labeling consistent with every named vertex constraint in {[}LW{]}, the
outer boundary has cyclic order

\[(8,3,21,23,11,13,6,2).\]

This order is a candidate used by the certificate, not a numerically
trusted hull. Its eight edges are tested against all 24 outer vertices
and all 24 reflected plug vertices. Each of these 384 determinants has
the affine form \(A+sD\) over \(\mathbb Q(t)\). Nonnegative values at
\(s=0\) and \(s=2-t\) establish every inequality throughout the interval
by linear interpolation. Exact zero tests are performed in
\(\mathbb Q[t]/(t^3-t^2-t-1)\); nonzero signs are enclosed using Arb.
Distinctness of the boundary vertices, noncollapsing edges, and strictly
positive cyclic turns are also checked throughout the interval,
establishing that the supporting polygon is the actual convex hull. A
second implementation of the field arithmetic recomputes all 384
coefficients.

The first failing test is the reflected vertex \(p=(-t,q,-1)\), labeled
20, against the oriented edge

\[a=(-t,-1,-q),\qquad b=(-1,-q,-t),\]

labeled \(8\to3\). Direct reduction of its determinant gives

\[\det\bigl(Q_s(b)-Q_s(a),Q'_s(p)-Q_s(a)\bigr)
=2(t-1)^2(2-t-s).\]

The positive prefactor \(2(t-1)^2\) makes the root exactly \(s_*=2-t\).
The half-turn \(\operatorname{diag}(-1,-1,1)\) gives the identical
obstruction at reflected vertex 17 against edge \(21\to23\). The first
edge remains supporting for \textbf{every} \(s\in[0,1/3]\), as its
outer-vertex affine inequalities are nonnegative at both endpoints. Its
plug determinant is strictly negative for \(s\in(2-t,1/3]\), proving the
converse.

Finally \(K_{u_\theta}K_0=R_z(2\theta)\), and
\(\pi_{u_\theta}K_{u_\theta}=\pi_{u_\theta}\). This proves the rotation
equivalence. The plug rotation preserves \(v_z\), and
\(a_4\perp u_\theta\), so both vertical support values are equal; the
contained shadow always touches. This observation also rules out strict
containment by translating this particular pair of shadows, since
translation cannot reduce their equal vertical widths.

At 384-bit working precision the certified enclosure is

\[\begin{aligned}
\theta_*\in[\,&0.1593506261845691446953081463393635724078,\\
&0.1593506261845691446953081463393635724079\,]\,\mathrm{rad}.
\end{aligned}\]

It is approximately \(9.13011834314267604063121785735^\circ\). The
decimal interval has exact rational endpoints and width \(10^{-40}\); it
is not a confidence interval or a sample-based estimate.

\newpage

\section{Comparison with the two
sources}\label{comparison-with-the-two-sources}

In {[}LW, Theorem 4.2{]}, \(\theta_0=\arctan(0.16071324478583)\) and
\(\theta_*\in[\theta_0,\theta_0+10^{-14}]\). Our check gives
\(\theta_*-\theta_0\approx8.64418\times10^{-15}\), within that bracket.
The preprint already establishes the arc and characterizes its endpoint
by the first signed-area zero. Our contribution is its simplification,
an independent containment certificate, and the explicit comparison
below. To our knowledge, the closed form is not stated in the supplied
October 3 version of {[}LW{]} or the inspected repository version of
{[}Tom7{]}. This is a claim about those versions, not a literature-wide
priority claim.

Tom 7's exact frame is

\[x=(1,0,t^2),\quad e_1=(-t^2,0,1),\quad Z=e_y.\]

The wall in question is the \(XZ\)-plane of outer views. On side 1 of
the cap chart, the tie has

\[u_T=x-\mu e_1,\qquad w=\mu e_y,\qquad
\phi=\arctan\mu,\qquad R_T=R_y(2\phi).\]

Because \(x\perp e_1\) and \(|x|=|e_1|\), \(\phi\) is the angle past
\(u^*=x/|x|\), namely the angle denoted \(\varepsilon\) in the markdown
explanation. With \(H_n=2nn^{\mathsf T}/|n|^2-I=-K_n\), the half-turn
convention used by the repository, \(H_{u_T}H_x=K_{u_T}K_x=R_y(2\phi)\).
Individually \(H_u\) negates a screen projection; \(K_u\) fixes it.
Their products coincide. The inner viewing direction is
\(R_T^{\mathsf T}u_T\), explaining the opposite-angle inner view in the
C++ comment.

The proper signed permutation \(S(X,Y,Z)=(-X,Z,Y)\), followed by
division by \(t\), maps Tom's solid to \(P\). If \(k=1-q>0\), then

\[Sx=C/k,\qquad Se_1=B/k,\qquad Se_y=a_4.\]

Consequently \(\phi=\varepsilon_{\rm Tom}=\theta\), \(\mu=\tan\theta\),
and the L, W chart variable is \(\varepsilon_{\rm LW}=-\mu\). These are
the \textbf{same family}. Directly in Tom's unnormalized screen
coordinates \(((e_1+\mu x)\cdot v,v_Y)\), the critical edge vertices and
plug vertex are

\[(t^2,-1,-t),\quad(t,-t^2,-1),\quad(t^2,-t,1),\]

and their determinant is

\[2(t-1)-2t(t+1)\mu=2t(t+1)(2-t-\mu).\]

Thus \(\mu_*=2-t=0.160713244785838867448\ldots\) and
\(\varepsilon_{\rm Tom,*}=\theta_*\); the relative rotation ends at
\(2\theta_*=0.318701252369138\ldots\) radians. The markdown states about
\(0.165\) \textbf{radians} (\texttt{snub-cube.md:57–60}), while the C++
comment states \(\mu\) up to about \(0.165\)
(\texttt{search/sc/sccapcert.cc:23–26}). The new certificate proves
failure at both \(\mu=0.165\) and \(\varepsilon_{\rm Tom}=0.165\)
radians: the critical signed determinants are respectively enclosed near
\(-0.0447731222\) and \(-0.0605847034\). Hence changing between angle
and Cayley parameter does not resolve either approximate bound. This
corrects a descriptive range, without showing any failure of the global
cap/tube exclusion certificates. Their \(\mu\leq1/3\) domain need not
consist entirely of touching poses.

All 503 tracked files of {[}Tom7{]} were enumerated and all 500 text
files searched, including Lean, C++, generated data and Markdown. The
three binary images were excluded from textual search. The scan did not
cover repository history, external \texttt{nopert229} prototypes, linked
videos/papers, or unprovided {[}LW{]} supplements; \texttt{novelty.md}
records the evidence and limits.

\newpage

\section{Two remarks on {[}LW{]}}\label{two-remarks-on-lw}

\textbf{Weight signs in Lemma 5.2.} A common sign does not imply that
the weights are nonnegative. Flipping three negative weights also flips
their weighted support gap, so it requires rechecking the strict
inequality. An explicit counterexample to the general algebraic
criterion is

\[P_b=[-1,1]\times[-2,2]\times[-3,3],\quad u=e_1,\quad
R(x,y,z)=(z,x,y),\quad r=(1,1,1),\quad D=4.\]

In the \((y,z)\) screen the hole is \([-2,2]\times[-3,3]\) and the plug
is \([-1,1]\times[-2,2]\), strictly contained at translation zero.
Choose \(c=e_1,a_4=e_3,b=-e_2\), \(\varepsilon=\delta=0\), and these
tests:

{\def\LTcaptype{none} % do not increment counter
\begin{longtable}[]{@{}
  >{\raggedright\arraybackslash}p{(\linewidth - 12\tabcolsep) * \real{0.1304}}
  >{\raggedright\arraybackslash}p{(\linewidth - 12\tabcolsep) * \real{0.1304}}
  >{\raggedright\arraybackslash}p{(\linewidth - 12\tabcolsep) * \real{0.1304}}
  >{\raggedright\arraybackslash}p{(\linewidth - 12\tabcolsep) * \real{0.1304}}
  >{\raggedright\arraybackslash}p{(\linewidth - 12\tabcolsep) * \real{0.1304}}
  >{\raggedleft\arraybackslash}p{(\linewidth - 12\tabcolsep) * \real{0.1739}}
  >{\raggedleft\arraybackslash}p{(\linewidth - 12\tabcolsep) * \real{0.1739}}@{}}
\toprule\noalign{}
\begin{minipage}[b]{\linewidth}\raggedright
\(k\)
\end{minipage} & \begin{minipage}[b]{\linewidth}\raggedright
\(\nu_k\)
\end{minipage} & \begin{minipage}[b]{\linewidth}\raggedright
\(u\times\nu_k\)
\end{minipage} & \begin{minipage}[b]{\linewidth}\raggedright
plug vertex \(p_k\)
\end{minipage} & \begin{minipage}[b]{\linewidth}\raggedright
hole support vertex \(q_k\)
\end{minipage} & \begin{minipage}[b]{\linewidth}\raggedleft
\(\Delta_k\)
\end{minipage} & \begin{minipage}[b]{\linewidth}\raggedleft
\(g_k\)
\end{minipage} \\
\midrule\noalign{}
\endhead
\bottomrule\noalign{}
\endlastfoot
1 & \((0,0,-1)\) & \((0,1,0)\) & \((1,2,3)\) & \((1,2,3)\) & −1 & −4 \\
2 & \((0,1,0)\) & \((0,0,1)\) & \((1,2,3)\) & \((1,2,3)\) & −1 & −4 \\
3 & \((0,-1,1)\) & \((0,-1,-1)\) & \((-1,-2,3)\) & \((1,-2,-3)\) & −1 &
−8 \\
\end{longtable}
}

Here \(g_k=\langle N(r)p_k-Dq_k,u\times\nu_k\rangle\),
\(\sum\Delta_k(u\times\nu_k)=0\), and \(\sum\Delta_kg_k=16>0\). Each
hole support test is valid. The constant chart \(\Phi=(0,0,1,1,1)\) on
\([1,2]^5\), with distance variable \(y_0\) and \(m_F=1\), also meets
the divisibility conditions. Thus the common-sign inference and its
corresponding general criterion fail. This box is not a snub cube and is
not a counterexample to either snub solid's non-Rupert conclusion. The
needed correction is to require nonnegative weights and evaluate
condition (3) using those same weights. Whether the actual {[}LW{]}
certificates already do so is outside this check.

\textbf{Appendix A.4 and binary64 constants.} Section 5.3 interprets
decimal constants as binary64 values. However

\[\operatorname{fl}(1.2)=\frac{5404319552844595}{4503599627370496}<\frac65,
\quad \frac65-\operatorname{fl}(1.2)=\frac1{22517998136852480}.\]

For the special pair with \(s_2=1/5\), \(f_2=1,d=-1\) gives
\(|f_2-ds_2|=6/5\), violating the displayed bound if read as
\(\leq\operatorname{fl}(1.2)\). For example \(\mu=1/100\),
\(\varepsilon=\delta=-\mu\), \(r=(0,\mu,\mu)\) is a corresponding CORNER
point. The exact fraction \(6/5\) repairs this wording. It is already
admitted by the relevant exact chart bound, and the enlarged Shigh bound
\(5/4\) is larger. This discrepancy is 0.2 ulp at 1.2; it does not
introduce a new hole in the continuum cover checked by our audit.

\newpage

\section{Reproducibility, scope and author statement}
\label{reproducibility-scope-and-author-statement}

The public repository contains \texttt{code/} (six mathematical scripts),
\texttt{data/} (reference JSON certificates and provenance), and
\texttt{note/} (this note and the source-review record).
From the repository root run:
\begin{verbatim}
python -B self_test.py
python -B run_all.py
\end{verbatim}
Both commands execute the checks sequentially in one Python process with
numerical thread limits set to 1. The self-test verifies SHA-256 integrity,
reruns the mathematics and compares eight JSON certificates, and checks the
five-page PDF. Generated results go to \texttt{out/}. No upstream PDF,
clone, Lean build, C++ search or network access is needed. Historical input
hashes identify the reviewed sources; source-document inspection is a
separate procedure in the README and \texttt{novelty.md}.

The recorded environment is Python 3.13.15, NumPy 2.5.2, mpmath 1.3.0,
SymPy 1.14.0, python-flint 0.9.0, and pypdf 6.20.1. NumPy is recorded for
environment parity; the certificate path does not use it. mpmath is a
SymPy dependency; pypdf is used by the release self-test. Script hashes,
measured runtime and claim-to-command mapping are in the README,
\texttt{SHA256SUMS}, and \texttt{data/reproducibility.json}.

The reviewed preprint is DOI 10.5281/zenodo.23118432. The noperts snapshot
is commit \texttt{0824ef59\allowbreak{}8e863748\allowbreak{}1848092f\allowbreak{}f283d815\allowbreak{}cf0c04e8}.
The source search concerns that snapshot only, not repository history or
literature-wide priority. The check covers Theorem 4.2 and the local family
through $\mu=1/3$, plus two textual remarks. It does not certify all
poses/translations, either global non-Rupert result, or a global Nieuwland
constant.

Data and code availability. The version 1.0.1 release is available at
\url{https://github.com/ysmin1575/snub-cube-touching-arc}
and archived on Zenodo; its version-specific code DOI is
\url{https://doi.org/10.5281/zenodo.23397095}.

Use of AI. The author chose to audit the snub-cube preprint, directed the project and decided at each stage what to check and what to publish. AI systems were used extensively as research assistants: Claude (Anthropic) for planning, literature search, writing and running code, and drafting; GPT-6.1 Sol (OpenAI, via Codex) for independent derivations, the verification scripts and the first draft of this note. Every mathematical claim in this note is checked by the published scripts in exact or outward-rounded interval arithmetic, so the reader does not need to trust any AI system: running self\_test.py reproduces all certified statements.

\small


\textbf{References}

\begin{itemize}
\tightlist
\item
  \textbf{{[}LW{]}} L, W (creator name as registered in Zenodo).
  \emph{The Snub Cube and Snub Dodecahedron Are Not Rupert}. Preprint
  dated 3 October 2026; the PDF title includes ``the'' before ``Snub
  Dodecahedron.'' \href{https://doi.org/10.5281/zenodo.23118432}{DOI:
  10.5281/zenodo.23118432}, especially Theorem 4.2, Lemma 5.2 and
  Appendix A.4.
\item
  \textbf{{[}Tom7{]}} Tom 7. \emph{noperts}, snub-cube description and
  sources, inspected at commit
  \texttt{0824ef59\allowbreak{}8e863748\allowbreak{}1848092f\allowbreak{}f283d815\allowbreak{}cf0c04e8}, 14 October 2026;
  the snub-cube status entry reports verification on 13 October.
  \href{https://github.com/tom7/noperts/blob/0824ef598e8637481848092ff283d815cf0c04e8/snub-cube.md\#L53-L60}{Pinned
  description},
  \href{https://github.com/tom7/noperts/blob/0824ef598e8637481848092ff283d815cf0c04e8/search/sc/sccapcert.cc\#L6-L26}{C++
  chart description},
  \href{https://github.com/tom7/noperts/blob/0824ef598e8637481848092ff283d815cf0c04e8/proof/SnubCube/Caps.lean\#L35-L64}{exact
  Lean frame}.
\end{itemize}

\end{document}
